% English translation of checked_pages/esq_119_024.tex; source: cote 119, site 24, PDF page 25, typescript page 22.
% AI-translated by GPT-6 Astra (OpenAI, ChatGPT Pro mode).
\scanpage{119/24}
\mspage{22}

regular polyhedron under consideration is finite. In the case of the
base field $\mathbb{R}$ (or $\mathbb{C}$, which amounts to the same thing), and for $n = 2$, the finite
cases have been well known since antiquity - which does not exclude
the scheme-theoretic viewpoint's revealing new charms
in them; one can say, however, that by specializing the icosahedron
(for example) over finite base fields of arbitrary
characteristic, one always obtains an icosahedron, with its own
combinatorics and the same simple automorphism group of order 60.
The same remark applies to finite regular polyhedra
of higher dimension, studied systematically in two
fine books by Coxeter. The situation is quite different on starting
with an \uline{infinite} regular polyhedron, over a field such as $\mathbb{Q}$, say,
and ``specializing'' it over the prime fields $\mathbb{F}_p$ (an operation
well defined for every $p$ except finitely many primes).
It is clear that every regular polyhedron over a finite field is finite -
\uline{one thus finds infinitely many finite regular polyhedra for} $p$
\uline{varying}, \uline{whose combinatorial type}, or, what amounts to the same thing,
whose automorphism group, \uline{varies} ``\uline{arithmetically}'' with p.
This situation is particularly intriguing in the case
$n = 2$, where one has the relation made explicit in the preceding
paragraph between combinatorial 2-maps and algebraic curves
defined over number fields. In this case an infinite regular
polyhedron defined over any infinite field (and hence
over a sub-$\mathbb{Z}$-algebra with two generators \unsure{.}of it)\tnote{TN-024-01}{The small dot before the French ``de'' is unresolved; it remains at the corresponding ``of'' in the translation.}
thus gives rise to infinitely many algebraic curves defined over
number fields, which are Galois coverings of the standard projective
line ramified only at 0, 1, $\infty$. The optimal case is of course
that in which one starts with the universal regular 2-polyhedron,
or rather with the one deduced from it by passing to the fraction
field $\mathbb{Q}(\alpha_1,\alpha_2)$ of its base ring. This raises a host of
new questions, vague as well as precise, none of which
I have yet had the leisure to examine more closely - I shall
mention only this one: what exactly are the finite regular
2-maps, or, what amounts to the same thing, \unsure{.}the quotient groups\tnote{TN-024-02}{The unresolved dot before the French ``les'' is retained before the corresponding English ``the''.}

\margin{24}

